\subsection{借助延拓定义测度}

设 X 是一个局部紧空间；若 V 是 \(\mathcal{K}(\mathrm{X};\mathbf{C})\) 的一个稠密线性子空间，则显然 X 上在 V 上取值相同的两个测度 \(\mu_{1}, \mu_{2}\) 相等，且 V 上对 \(\mathcal{K}(\mathrm{X};\mathbf{C})\) 的拓扑所诱导的拓扑连续的每个线性形式可以（以唯一方式）延拓成 X 上的一个测度。对正测度而言，一个方便的判别法如下：

命题 9. --- 设 V 是 \(\mathcal{K}(\mathbf{X};\mathbf{R})\) 的一个线性子空间，具有下述性质：

\begin{enumerate}
\def\labelenumi{(\Alph{enumi})}
\setcounter{enumi}{15}
\tightlist
\item
  对每个紧子集 \(\mathbf{K}\) （\(\mathbf{X}\) 的），存在一个函数 \(f \in \mathbf{V}\) ，使得 \(f \geqslant 0\) 且 \(f(x) > 0\) 对所有 \(x \in \mathbf{K}\) 成立。
\end{enumerate}

在这些条件下，V 上对由 \(\mathcal{K}(\mathrm{X};\mathbf{R})\) 的序所诱导的序（Ch. II, §2, No. 1, 定义 1）而言的每个正线性形式可以延拓成 X 上的一个正测度（当 V 在 \(\mathcal{K}(\mathrm{X};\mathbf{R})\) 中稠密时它是唯一的）。

对每个函数 \(f \in \mathcal{K}(\mathbf{X};\mathbf{R})\) （支集为 \(K\) 者），存在一个函数 \(g \in V\) 使 \(f \leqslant g\) ；因为存在一个函数 \(h \geqslant 0\) （\(V\) 中的），使得 \(h(x) > 0\) 对所有 \(x \in K\) 成立；令 \(\alpha = \inf_{x \in K} h(x)\) ，于是 \(\alpha > 0\) ，且函数 \(g = (\alpha^{-1} \|f\|)h\) 满足要求。于是只须应用 No. 5 的 Th. 1 与 TVS, II, §3, No. 1 的命题 1。

